\section{Architettura generale del sistema}
\label{sec:architecture}

\subsection{Componenti e confini di responsabilità}

Il sistema adotta un'architettura distribuita con controllo centralizzato nel CUS, responsabile della gestione dello stato globale, delle modalità operative e della determinazione dell'apertura della valvola.
Le responsabilità dei diversi sottosistemi sono illustrate nella Figura~\ref{fig:architecture}.

\begin{figure}[htbp]
\centering
\begin{tikzpicture}[node distance=1.2cm and 1.4cm]
\node[block] (tms) {\textbf{TMS}\\ESP32\\Sonar e LED};
\node[block,right=of tms] (broker) {\textbf{Broker MQTT}};
\node[block,below=of broker] (cus) {\textbf{CUS}\\PC/Java\\Controllo e stato};
\node[block,left=of cus] (wcs) {\textbf{WCS}\\Arduino UNO\\Servo e LCD\\Comandi locali};
\node[block,right=of cus] (dbs) {\textbf{DBS}\\Browser\\Supervisione e comandi};
\draw[flow] (tms) -- node[above,font=\footnotesize]{misure} (broker);
\draw[flow] (broker) -- node[right,font=\footnotesize]{MQTT} (cus);
\draw[flow,<->] (wcs) -- node[above,font=\footnotesize]{seriale} (cus);
\draw[flow,<->] (cus) -- node[above,font=\footnotesize]{HTTP} (dbs);
\end{tikzpicture}

\caption{Architettura del sistema e collegamenti tra i sottosistemi.}
\label{fig:architecture}
\end{figure}

Il TMS acquisisce il livello dell'acqua tramite un sensore sonar e trasmette le misure al CUS.
Quest'ultimo elabora i dati ricevuti, applica la politica di controllo e rende disponibile lo stato del sistema.
Il WCS gestisce l'attuatore e l'interfaccia locale, consentendo all'operatore di richiedere il cambio di modalità tramite pulsante e la regolazione dell'apertura tramite potenziometro.
Il DBS permette il monitoraggio e il controllo remoto attraverso un'interfaccia web.

\subsection{Comunicazione e propagazione dello stato}

La comunicazione tra TMS e CUS avviene tramite MQTT, utilizzando un broker per lo scambio delle misure.
Il WCS comunica con il CUS attraverso una connessione seriale, inviando le richieste dell'operatore e ricevendo gli aggiornamenti di stato necessari al controllo dell'attuatore e del display.

Il DBS comunica con il CUS tramite HTTP, recuperando periodicamente lo stato del sistema e lo storico delle misure.
Attraverso la stessa interfaccia può inviare richieste di cambio modalità e regolazione dell'apertura.

I dettagli dei protocolli e dei messaggi sono descritti nella Sezione~\ref{subsec:integration-contracts}.

\subsection{Stato globale e disponibilità}

Il CUS gestisce le modalità operative \code{AUTOMATIC} e \code{MANUAL} e monitora la disponibilità delle misure provenienti dal TMS.
In assenza di misure valide per un intervallo di tempo prestabilito, il sistema passa allo stato \code{UNCONNECTED}.

Il DBS distingue questa condizione dalla propria indisponibilità di comunicazione con il CUS, rappresentata come \code{NOT AVAILABLE}.
Il backend può quindi essere raggiungibile anche quando le misure del TMS non sono disponibili.

\subsection{Scelte architetturali}

La politica di controllo è gestita esclusivamente dal CUS.
Le interfacce locale e remota inviano richieste che vengono valutate in base allo stato corrente del sistema, senza una priorità prestabilita tra i due operatori.

Nel backend, la logica di controllo è separata dai componenti responsabili delle comunicazioni MQTT, seriali e HTTP.
L'accesso allo stato condiviso è sincronizzato per gestire le operazioni concorrenti, mentre le operazioni di comunicazione vengono eseguite al di fuori delle sezioni critiche.